% Standalone uncompiled research note. No external TeX inputs are required.
% Preserve this draft as an intermediate backup.
\documentclass[11pt]{article}
\usepackage[T1]{fontenc}
\usepackage[utf8]{inputenc}
\usepackage[margin=1in]{geometry}
\usepackage{amsmath,amssymb,amsthm,mathtools,enumitem}
\usepackage[hidelinks]{hyperref}
\newtheorem{theorem}{Theorem}[section]
\newtheorem{proposition}[theorem]{Proposition}
\newtheorem{lemma}[theorem]{Lemma}
\newtheorem{corollary}[theorem]{Corollary}
\theoremstyle{remark}\newtheorem{remark}[theorem]{Remark}
\newcommand{\R}{\mathbb R}
\newcommand{\Id}{\mathrm I}
\newcommand{\rank}{\operatorname{rank}}
\newcommand{\tr}{\operatorname{tr}}
\newcommand{\diag}{\operatorname{diag}}
\newcommand{\op}{\mathrm{op}}
\title{Controlled training responses in two layer matrix models}
\author{}
\date{5 November 2026}
\begin{document}
\maketitle
\begin{abstract}
For two-factor real matrix models, we classify equality of full output
responses under changing squared-loss targets. Except when both external
dimensions are one, simultaneous gradient descent at any fixed nonzero
step size distinguishes exactly the full block Gram state; target words
of length at most two suffice. Controlled gradient flow retains a scalar
Gram-shift ambiguity. At a fixed hidden width the distinction exists
exactly above an explicit width threshold. A concrete pair is
indistinguishable under every stationary target but distinguishable after
one target change. Its optimal two-update separation under bounded targets
is cubic in the step size, giving a sharp adversarial-noise boundary for
restarted experiments of that depth. The proof uses established covariance,
mobility and conservation-defect identities. This is a self-contained
corollary-level refinement with precise experimental quantifiers, not a
claim of historical novelty, a new optimization theory, or empirical benefit.
\end{abstract}

\section{The model and the experiment}
Fix $m,n\geq1$. A state consists of real matrices
$A\in\R^{m\times h}$ and $B\in\R^{h\times n}$ at a finite positive hidden
width $h$. Its observed output is the \emph{whole} matrix $F=AB$.
Each target $Y\in\R^{m\times n}$ specifies the loss
$\tfrac12\|AB-Y\|_F^2$ and the simultaneous update
\begin{equation}\label{eq:factor-step}
 A^+=A-\eta EB^{\mathsf T},\qquad
 B^+=B-\eta A^{\mathsf T}E,\qquad E=F-Y,
\end{equation}
where the known step size $\eta\ne0$ is fixed and the same in both layers.
Ordinary descent uses $\eta>0$; the algebra below holds for either sign.
Both gradients use the old factors. No bias, momentum, weight decay,
projection or alternating update is included.

Write $F_0$ for the initial output and $F_k$ for the output after $k$
updates. Two states, possibly with different widths, are \emph{response
equivalent} if they give the same $F_k$ for every finite common target
word $(Y_0,\ldots,Y_{k-1})$, including $k=0$. The phrase ``through two
updates'' means equality of $F_0,F_1,F_2$ under every target word of length
at most two. It does not mean that one pair of probes identifies every
unknown state. The target can change within a word. Targets and output
matrices are unrestricted unless a later statement imposes a norm bound.

These are exact-model comparisons. Quantifying over all hypothetical
experiments does not supply reset access, calibrated interventions,
noiseless readings, or evidence that an application follows
\eqref{eq:factor-step}.

\section{The closed Gram state}
Set
\begin{equation}\label{eq:gram}
 L=AA^{\mathsf T},\quad R=B^{\mathsf T}B,\quad
 C=\begin{pmatrix}A\\B^{\mathsf T}\end{pmatrix},\quad
 G=CC^{\mathsf T}=\begin{pmatrix}L&F\\F^{\mathsf T}&R\end{pmatrix},
 \quad J=\diag(\Id_m,-\Id_n).
\end{equation}
The realizable states at width $h$ are precisely $G\succeq0$ with
$\rank G\leq h$: a Gram factorization constructs $C$, and its two blocks
recover $A,B$. In particular, full positivity of $G$, not just positivity
of $L$ and $R$, is required.

Direct multiplication of \eqref{eq:factor-step} gives
\begin{align}
 F^+&=F-\eta(LE+ER)+\eta^2EF^{\mathsf T}E,\label{eq:Fstep}\\
 L^+&=L-\eta(EF^{\mathsf T}+FE^{\mathsf T})+\eta^2ERE^{\mathsf T},
 \label{eq:Lstep}\\
 R^+&=R-\eta(E^{\mathsf T}F+F^{\mathsf T}E)+\eta^2E^{\mathsf T}LE.
 \label{eq:Rstep}
\end{align}
Equivalently, for
$S_E=\left(\begin{smallmatrix}0&E\\E^{\mathsf T}&0\end{smallmatrix}\right)$,
\begin{equation}\label{eq:congruence}
 C^+=(\Id-\eta S_E)C,\qquad
 G^+=(\Id-\eta S_E)G(\Id-\eta S_E).
\end{equation}
These identities need no balancedness, full-rank or genericity assumption.
They include $F=0$ and singular Gram matrices. They are established
covariance dynamics; Section~\ref{sec:prior} specifies the close sources.

\begin{lemma}[Exactly what one update reveals]\label{lem:one}
Two states have equal initial and one-update responses for every target
if and only if $G_2-G_1=cJ$ for some $c\in\R$.
\end{lemma}
\begin{proof}
Initial equality fixes $F$. As $Y$ varies, $E=F-Y$ is arbitrary.
Equation~\eqref{eq:Fstep} therefore gives
$\Delta L E+E\Delta R=0$ for every $E$, where $\Delta$ means state 2
minus state 1. For $E=e_i e_j^{\mathsf T}$ the $(k,\ell)$ entry is
$\Delta L_{ki}\delta_{j\ell}+\delta_{ki}\Delta R_{j\ell}$.
Taking $k\ne i,\ell=j$ and then $k=i,\ell\ne j$ kills all off-diagonal
entries. Taking $k=i,\ell=j$ shows
$\Delta L_{ii}=-\Delta R_{jj}$ for every $i,j$. Hence
$\Delta L=c\Id_m$, $\Delta R=-c\Id_n$. Conversely these differences
cancel in \eqref{eq:Fstep}.
\end{proof}

\section{Two updates and the scalar exception}
\begin{theorem}[Complete discrete response classification]\label{thm:discrete}
If $mn>1$, the following conditions are equivalent:
\begin{enumerate}[label=(\roman*)]
\item the two states are response equivalent;
\item their responses agree through two updates;
\item their initial Gram matrices $G$ agree.
\end{enumerate}
If $m=n=1$, response equivalence is instead equality of
$(f,s)=(F,L+R)$ and is already determined by initial and one-update
responses.
\end{theorem}
\begin{proof}
For $mn>1$, (i) implies (ii). By Lemma~\ref{lem:one}, (ii) gives an
initial difference $cJ$. After any first residual $E$, the outputs agree
and the other differences are
\[
 \Delta L_1=c(\Id_m-\eta^2EE^{\mathsf T}),\qquad
 \Delta R_1=-c(\Id_n-\eta^2E^{\mathsf T}E).
\]
For any second residual $Z$, choose the common second target $Y_1=F_1-Z$.
The terms depending only on $F_1$ cancel in \eqref{eq:Fstep}, leaving
\begin{equation}\label{eq:cubic}
 \Delta F_2=\eta^3c\bigl(EE^{\mathsf T}Z-ZE^{\mathsf T}E\bigr).
\end{equation}
If $n>1$, choose the rectangular elementary matrices
$E=e_1e_1^{\mathsf T}$ and $Z=e_1e_2^{\mathsf T}$; the parenthesis is $Z$.
If $n=1,m>1$, choose $E=e_1,Z=e_2$; it is $-Z$.
Since $\eta\ne0$, equality forces $c=0$, proving (iii).
The closed equations \eqref{eq:Fstep}--\eqref{eq:Rstep} prove
(iii)$\Rightarrow$(i) by induction.

When $m=n=1$, write $e=f-y$. The equations reduce to
\begin{equation}\label{eq:scalar}
 f^+=(1+\eta^2e^2)f-\eta es,\qquad
 s^+=(1+\eta^2e^2)s-4\eta ef.
\end{equation}
Thus $(f,s)$ is sufficient for every future response, and any first
residual $e\ne0$ recovers $s$ from $f,f^+$. The cubic parenthesis is
identically zero in this case.
\end{proof}

\begin{remark}[The target quantifier]
The separating word in the proof is a legitimate common \emph{fixed}
word once the two candidate states and $E$ have been fixed: their shared
$F_1$ determines the same matrix $Y_1$. No adaptive policy is needed for
the equivalence theorem. A procedure can instead choose $Y_1$ after
observing $F_1$, but then exact access or an appropriate error analysis
is an additional assumption. A bound on distinguishing-word depth is
not a bound of two measurements for reconstructing every unknown Gram
matrix.
\end{remark}

\section{Controlled flow and global existence}
Consider the actual controlled gradient flow
\begin{equation}\label{eq:flow}
 \dot A=-EB^{\mathsf T},\qquad \dot B=-A^{\mathsf T}E,
 \qquad E=AB-Y(t),
\end{equation}
with a common locally bounded piecewise-continuous target $Y(t)$.
Flow equivalence means equality of the entire output path for every
such target signal on every finite time interval.

\begin{theorem}[Complete controlled-flow classification]\label{thm:flow}
Every initial state has a unique global forward solution to
\eqref{eq:flow} for the stated controls. Two states are flow equivalent
if and only if $G_2-G_1=cJ$ for some $c\in\R$.
\end{theorem}
\begin{proof}
The Gram equations are
\begin{align*}
 \dot F&=-LE-ER,\\
 \dot L&=-EF^{\mathsf T}-FE^{\mathsf T},&
 \dot R&=-E^{\mathsf T}F-F^{\mathsf T}E.
\end{align*}
Local existence and uniqueness hold piecewise, or in the usual
Carath\'eodory formulation. For
$V=(\|A\|_F^2+\|B\|_F^2)/2$, almost everywhere,
\[
 \dot V=-2\|F\|_F^2+2\langle F,Y\rangle_F
          \leq \tfrac12\|Y\|_F^2.
\]
Hence $V$ is bounded on every finite interval. The locally Lipschitz
parameter system cannot have finite-time blowup, giving global existence.

Equal initial output and equal initial derivatives for every constant
target imply $\Delta LE+E\Delta R=0$ for every $E$.
Lemma~\ref{lem:one}'s matrix argument gives necessity.
For sufficiency, the shift $\Delta L=c\Id_m$, $\Delta R=-c\Id_n$
leaves $\dot F$ unchanged, while the other two equations preserve this
constant difference whenever $F$ agrees. Uniqueness of the closed system
therefore gives equal outputs globally. Equivalently,
$\dot G=-S_EG-GS_E$ and $S_EJ+JS_E=0$.
\end{proof}

Flow equivalence is a feasible Gram-shift relation, not the hidden
orthogonal gauge: a nonzero shift changes $G$ and can change its rank.
Theorems~\ref{thm:discrete} and \ref{thm:flow} compare exact equivalence
relations. They do not contradict convergence of small-step numerical
trajectories: the separating signal in \eqref{eq:cubic} tends to zero.
If finite gradient descent is the admitted operation, the difference is
an actual response of that operation. If Euler is only an approximation
to a physical flow, its extra distinction is an approximation artifact,
not evidence that the physical states are distinguishable.

\section{A concrete pair and the need to change targets}
For $m=1,n=2$, take
\[
 P:\quad A=(1),\ B=(0\ \ 0),\qquad
 Q:\quad A=(0\ \ 0),\ B=\Id_2.
\]
Both have $F_0=0$, while
$G_P=\diag(1,0,0)$ and $G_Q=\diag(0,1,1)$.
Thus $G_Q-G_P=-J$, and every controlled flow gives identical outputs.
For the two discrete targets
\[
 Y_0=(-1,0),\qquad Y_1=(-\eta,-1),
\]
straight substitution gives
\begin{align*}
 F_1(P)=F_1(Q)&=(-\eta,0),\\
 F_2(P)&=(-\eta,-\eta),&
 F_2(Q)&=(-\eta,-\eta-\eta^3).
\end{align*}
The same example has a common installed width two by padding $P$ to
$A=(1\ \ 0)$, $B=0_{2\times2}$.

\begin{proposition}[Stationary targets do not distinguish this pair]
For each single target $Y\in\R^{1\times2}$, held fixed throughout a run,
$P$ and $Q$ have identical discrete output sequences for every finite
number of updates. This holds for every nonzero step size.
\end{proposition}
\begin{proof}
For $Y=0$ both states remain at zero output. Otherwise an orthogonal
input-coordinate change carries $Y$ to $(y,0)$, $y>0$. In $Q$ a hidden
orthogonal change restores $B=\Id_2$; it does not change any output.
Along the first coordinate, $P$ starts with scalar factors $(a,b)=(1,0)$
and $Q$ with $(a,b)=(0,1)$. The perpendicular column in $Q$ remains
orthogonal to the active hidden coordinate and contributes zero output.
The active factors in both states follow the scalar update with target
$y$. Swapping them preserves their output and their update law, or use
\eqref{eq:scalar} with common $(f,s)=(0,1)$. Induction proves the result.
\end{proof}
Thus even trying arbitrarily many stationary targets on reset copies
misses this pair. The successful depth-two word changes the target within
one run. The larger pair
$\diag(\Id_m,0_n)$, $\diag(0_m,\Id_n)$ has the same stationary-target
property: an SVD of a fixed target splits it into swapped scalar systems
and unexcited directions.

\section{Width, quotients and finite access}
\begin{proposition}[The exact fixed-width boundary]\label{prop:width}
At a common fixed hidden width $h$, a nonzero feasible Gram shift can
occur if and only if $h\geq\max(m,n)$. Consequently, for $mn>1$:
\begin{enumerate}[label=(\roman*)]
\item if $h<\max(m,n)$, initial and one-update response equality already
forces full response equivalence;
\item if $h\geq\max(m,n)$, the two-update bound is sharp and controlled
flow is strictly less discriminating than discrete training.
\end{enumerate}
For $m=n=1$ the two regimes have the same $(f,s)$ quotient at every width.
\end{proposition}
\begin{proof}
If $G_2-G_1=cJ$ with $c>0$, then
$L_2=L_1+c\Id_m\succ0$ and $R_1=R_2+c\Id_n\succ0$.
Their ranks force $h\geq m,n$. Reverse the states for $c<0$.
Conversely $\diag(\Id_m,0_n)$ and $\diag(0_m,\Id_n)$ have ranks
$m,n$, admit width-$h$ realizations when $h\geq\max(m,n)$, and differ
by $-J$. Lemma~\ref{lem:one} and Theorem~\ref{thm:discrete} then give
both assertions and sharpness.
\end{proof}

At fixed width, two factors have the same $G$ exactly when they differ
by a hidden orthogonal transformation
$A\mapsto AU$, $B\mapsto U^{\mathsf T}B$.
Indeed equal Gram matrices give an isometry between the spans of the
rows of $C^{\mathsf T}$, which extends to an orthogonal map of the hidden
space. Across widths first pad by zero hidden coordinates. Thus, when
$mn>1$, the behavioral quotient agrees with this familiar Gram orbit
quotient. Its universal factorization property for invariant functions
is not a new universal construction.

Among all realizations of a nonscalar \emph{discrete} behavior, the minimum
hidden width is $\rank G$, allowing width zero; if widths must be positive
it is $\max(1,\rank G)$. Necessity is the rank bound; sufficiency is a Gram
factorization. Installed width is not identifiable because zero hidden
coordinates can always be appended. For flow the corresponding minimum
is instead
\[
 \min\{\rank(G+cJ):c\in\R,\ G+cJ\succeq0\},
\]
with the same positive-width convention. This target is constant on each
flow class. The example shows that flow observations cannot determine
the minimum width needed to reproduce all \emph{discrete} responses.
For scalar interfaces, feasibility gives $s\geq2|f|$ and every such
$(f,s)$ has a width-one realization: choose $a,b$ with
$a^2+b^2=s$, $ab=f$.

\begin{proposition}[A finite reset reconstruction]
For $mn>1$, exact full-matrix observations and reproducible copies of an
initial state permit recovery of $G$ using a finite family of one-update
experiments and one two-update experiment.
\end{proposition}
\begin{proof}
Observe $F$ and, for a chosen residual $E$, recover
\[
 K(E)=LE+ER=\frac{F-F^+(E)+\eta^2EF^{\mathsf T}E}{\eta}.
\]
Use $E=e_i e_1^{\mathsf T}$ for $1\leq i\leq m$ and
$E=e_1e_j^{\mathsf T}$ for $2\leq j\leq n$.
The entry formula in Lemma~\ref{lem:one} recovers all off-diagonal
entries of $L,R$ and their diagonal sums. It determines the representative
$L^0=L-a\Id_m$, $R^0=R+a\Id_n$ with $L^0_{11}=0$, where $a=L_{11}$.
This formal representative need not be positive semidefinite; it is
used only for polynomial calculation. Compare the actual second output
for the elementary $E,Z$ of Theorem~\ref{thm:discrete} with the output
predicted from $(L^0,F,R^0)$. Equation~\eqref{eq:cubic} recovers $a$
by division by $\eta^3$ in the selected entry, with its known sign.
\end{proof}
This procedure is not asserted to minimize the number of resets or
physical resources. Without reset or reproducible preparation, it does
not establish identification from one unknown training trajectory.

\section{A sharp bounded-depth noise boundary}
\begin{lemma}\label{lem:norm}
For rectangular $E,Z\in\R^{m\times n}$,
\[
 \|EE^{\mathsf T}Z-ZE^{\mathsf T}E\|_F
       \leq\|E\|_{\op}^2\|Z\|_F.
\]
\end{lemma}
\begin{proof}
An SVD $E=U\Sigma V^{\mathsf T}$ makes the $(i,j)$ entry of
$U^{\mathsf T}ZV$ acquire the multiplier $\lambda_i-\mu_j$, where
$\lambda_i,\mu_j$ are squared singular values, with zeros in the extra
dimensions. Both lie in $[0,\|E\|_{\op}^2]$. Sum the squared entry bounds.
\end{proof}

\begin{theorem}[Exact two-update discrimination threshold]\label{thm:noise}
For the concrete pair $P,Q$ above and target bound $M>0$,
\begin{equation}\label{eq:maxgap}
 \sup_{\|Y_0\|_F,\|Y_1\|_F\leq M}
     \|F_2(Q)-F_2(P)\|_F=|\eta|^3M^3.
\end{equation}
Suppose every output observation may have arbitrary additive error of
Frobenius norm at most $\varepsilon\geq0$, while targets, step size and
reset preparation are exact. Guaranteed discrimination of this declared
pair by reset experiments of depth at most two is possible if and only if
$|\eta|^3M^3>2\varepsilon$. Above the threshold one trial suffices;
at or below it no number of such trials, even adaptively chosen, guarantees
discrimination.
\end{theorem}
\begin{proof}
For this pair $F_0=0$, $E=-Y_0$ and the common $F_1=-\eta E$.
The next residual is $Z=-\eta E-Y_1$.
Since $EE^{\mathsf T}E=EE^{\mathsf T}E$, the term proportional to $E$
in \eqref{eq:cubic} cancels. With $c=-1$, the difference is
$\eta^3(EE^{\mathsf T}Y_1-Y_1E^{\mathsf T}E)$.
Lemma~\ref{lem:norm} gives the upper bound in \eqref{eq:maxgap};
$Y_0=-M(1,0)$, $Y_1=-M(0,1)$ attains it. These target norms are $M$
independently of $\eta$.

Two closed radius-$\varepsilon$ norm balls intersect exactly when their
centers have distance at most $2\varepsilon$. The initial and first
outputs of the pair coincide. Above the threshold the maximally separating
second-output balls are disjoint. At or below it an adversary can report
a common point in the intersecting second-output balls for every trial.
An adaptive protocol receives identical possible histories and therefore
chooses identical next targets and reset experiments. Repeating this
argument preserves ambiguity for any number of permitted trials.
\end{proof}
This is a worst-case bounded-error result for one candidate pair and
observation depth at most two. It is not an all-horizon impossibility:
longer changing-target runs can have different separation. It does not
apply unchanged to independent stochastic noise, which may average down,
or to uncertain target calibration. At fixed $M,\varepsilon>0$, exact
nonzero-step distinguishability therefore fails to imply warranted
identification in this experiment class as $\eta\to0$.

\section{What is inherited and what the comparison establishes}
\label{sec:prior}
The original project manuscript \emph{Learning generated characterizations}
\cite{ProjectLearning} treats training interventions as inputs and already
proves exact response closure and finite identification for scalar chains
of arbitrary depth. Its distinction between generated tests and acquired
observations remains in force. The present two-factor matrix model varies
hidden width and is not that scalar-chain theorem renamed. The common
operational construction and the changing-experiments note
\cite{ProjectCommon,ProjectChanging} supply the broader behavioral,
concept and supported-topology setting; no new FCA or topological result
is asserted here.

The closest algebraic antecedent is Zhang--Peyr\'e \cite{ZP}:
Proposition 2.1, equation (9), and Appendix C.5, equation (41) and
Lemma C.14, contain the covariance closure, discrete predictor update
and transported-indefinite-form defect used here, after adjusting
normalizations and block order. These are not new mechanisms.
Their stated results concern convergence and support margins; the
all-target two-initial-state classification is a different conclusion.

Tao \cite{Tao}, equations (3)--(4), Appendix B.2 and Section 7, uses the
same block Gram state, scalar-identity mobility ambiguity and constant
$J$-shift of continuous covariance dynamics. Lemma 3.8 treats effective
width reduction. The principal quantifier there concerns all autonomous
losses from one initialization. The present comparison is between two
initial states under target sequences; its stationary-target counterexample
makes the difference substantive. It should not be described as a new
continuous Gram-shift principle.

Marcotte--Peyr\'e--Gribonval \cite{MPG}, Theorems 3.8--3.9,
Proposition 3.10 and Appendix I, equations (26)--(27), give the close
intrinsic-dynamics, scalar and mobility-kernel antecedents.
Kunin et al.\ \cite{NeuralMechanics}, Sections 4--6 and Appendix I,
already establish the broad phenomenon that finite learning rates break
conservation laws of gradient flow. Pang et al.\ \cite{Pang}, Sections
3--5 and Appendix C, analyze early finite-step linear-network updates
and losses. Neither a Gram closure nor a two-step expansion alone is an
originality claim.

The narrow result is the proved response-equivalence classification,
its scalar and fixed-width boundary, the switched-target requirement,
and the stated rank and bounded-depth noise consequences. A targeted
primary-source audit did not locate that exact combination, but absence
from an inspected set is not priority evidence strong enough to certify
novelty. The proof is a short refinement of known identities. No new
optimizer, optimization rate, practical training advantage or general
universal property has been established.

\section{Verification and retained research files}
The universal statements above are justified by their proofs, not by
finite tests. The companion \texttt{check\_gram\_behavior.py} uses
SymPy exact arithmetic and checks 108 factor-versus-Gram identities,
108 cubic shift identities, 96 finite reconstructions, 12 scalar cases,
108 rank and congruence cases, symbolic separation and bounded-target
witnesses, eight scalar-ambiguity steps, and 30 stationary-target cases.
Its result is preserved in \texttt{check\_results.json}.

An independent \texttt{verify.py} additionally checks exact observation
ranks for all $1\leq m,n\leq3$, the symbolic concrete witness and scalar
closure, numerical update identities and exponential-flow shifts, and
120 stationary-target updates in four rectangular or square dimensions.
The largest reported numerical discrepancy is below $3\times10^{-15}$;
this is numerical regression evidence, not an exact proof.
Its proof, results and primary-source audit are retained alongside the
original working note, script and intermediate TeX draft. No formal
proof-assistant verification is claimed. This standalone TeX source has
been checked for source integrity and deliberately remains uncompiled.

\begin{thebibliography}{99}
\bibitem{ZP} S. Zhang and G. Peyr\'e,
\emph{Global Exponential Convergence of Two-Layer Linear Network Training},
arXiv:2610.09356v1, 7 October 2026.
\url{https://arxiv.org/html/2610.09356v1}.
\bibitem{Tao} S. Tao,
\emph{When Parameter Gradient Flows Close in Representation Space},
arXiv:2610.27474v1, 20 October 2026.
\url{https://arxiv.org/html/2610.27474v1}.
\bibitem{MPG} S. Marcotte, G. Peyr\'e and R. Gribonval,
\emph{Intrinsic training dynamics of deep neural networks}, ICLR 2026,
arXiv:2508.07370v2, 12 March 2026.
\url{https://arxiv.org/html/2508.07370v2}.
\bibitem{NeuralMechanics} D. Kunin, J. Sagastuy-Brena, S. Ganguli,
D. L. K. Yamins and H. Tanaka,
\emph{Neural Mechanics: Symmetry and Broken Conservation Laws in Deep
Learning Dynamics}, ICLR 2021, arXiv:2012.04728v2.
\url{https://arxiv.org/html/2012.04728}.
\bibitem{Pang} T. Pang, V. Kothapalli, S. Deng, H. Wang, D. Zhou and Y. Yang,
\emph{Balancing Learning Rates Across Layers: Exact Two-Step Dynamics and
Optimal Scaling in Linear Neural Networks}, arXiv:2606.00340v1,
29 May 2026. \url{https://arxiv.org/html/2606.00340v1}.
\bibitem{ProjectLearning} Previous project manuscript,
\emph{Learning generated characterizations}, 5 November 2026.
Recovered source \texttt{learning\_generated\_characterizations.tex}.
\bibitem{ProjectCommon} Previous project manuscript,
\emph{One operational construction for residual automata, formal concepts,
and topology}, 4 November 2026.
Recovered source \texttt{common\_generation\_note.tex}.
\bibitem{ProjectChanging} Previous project manuscript,
\emph{Changing experiments and composing characterizations},
5 November 2026. Source
\texttt{Changing\_experiments\_and\_composing\_characterizations.tex}.
\end{thebibliography}
\end{document}
